input implies $P(\zeta)\ne0$, so this integer is nonzero. This argument
uses reciprocity as a Laurent identity and does not assume that a leading
coefficient survives a finite-field specialization.
\end{proof}


The Fibonacci slice
\[
\Res(\Phi_K(X),X^2+3X+1)=F_K^2
\]
shows that individual collision resultants can be exponentially large.  At
$K=131$, $p=F_{131}$, a Lucas primality certificate and independent Python/C++
arbitrary-precision computations give exact nonzero Cayley excess $32$
(finite certificate:
\path{computations/verification/cayley_resultant_fibonacci.py} and
\path{computations/cpp/cayley_fibonacci_verifier.cpp}); this is
a finite certificate, not an asymptotic theorem.  The primitive-divisor input
for the associated rational model is from \cite{Yabuta2001}.

Let $S_K$ be the primitive squarefree nonzero-difference polynomial obtained
from the resultant of the characteristic-zero Cayley polynomial with its
translate.  For a split prime $p\equiv1\pmod K$, with $H$ and
$A=c(H\setminus\{1\})$ as in the preceding sections, define the
\emph{Cayley excess}
\begin{equation}\label{eq:Delta-def}
\Delta_K(p)
 =\sum_{d\ne0}r_{A-A}(d)^2-(K-1)(2K-5)
 =\sum_tC_H(t)^2-(K-1)(2K-5)\in\mathbf Z_{\ge0}.
\end{equation}
The two expressions agree by \eqref{eq:energyCH} and
$r_{A-A}(0)=|A|=K-1$.  Nonnegativity holds because the identity and central
inversion-swap collisions are distinct ordered collision quadruples already
contributing exactly $(K-1)(2K-5)$ to the first sum: $\Delta_K(p)$ counts
the remaining, non-forced, ordered collision quadruples.  Then for split
primes
\begin{equation}\label{eq:disc-class}
\Delta_K(p)>0\iff p\mid\Disc(S_K),
\qquad
\Delta_K(p)\le4v_p(\Disc(S_K)).
\end{equation}
The exact global form of $\Disc(S_K)$ comes from the ratio-orbit
factorization.  Write $n=K-1$ and $m=(K-1)/2$.  Then
$S_K=\prod_{[r]}F_r$, the product over the $m$ ratio classes
$[r]$, $r\sim r^{-1}$, where $F_r$ is the degree-$n$ Galois-orbit factor of
the differences $c(\zeta^a)-c(\zeta^{ra})$.  For one orbit, the index
formula in the maximal real subfield, whose discriminant is $K^{m-1}$,
gives $|\Disc(F_r)|=2^{n^2}K^{K-2}I_r^4$ for an integer $I_r>0$.  For two
distinct classes, $|\Res(F_r,F_s)|=2^{n^2}K^{K+1}M_{r,s}^2$; the square
comes from writing the even half-degree factors as $H_r(Y)=1+KP_r(Y)$ and
using $\deg(P_s-P_r)=m-1$.  That degree drop is the one place where
Dedekind sums enter: the coefficient of $Y^{m-1}$ in $P_r$ is
$(K-1)(K-2)/12-Ks(r,K)$, and by the prime-modulus equality criterion
\cite[Corollary~1.2]{JRW2011} (pp.~2197--2202), $s(r,K)=s(s,K)$ if and
only if $r\equiv s$ or $rs\equiv1\pmod K$, so distinct ratio classes have
distinct coefficients (for $K\ge5$; at $K=3$ there is a single ratio class
and the formula is a direct check).  Since
$\Disc(\prod_rF_r)=\prod_r\Disc(F_r)\prod_{r<s}\Res(F_r,F_s)^2$, collecting
the powers of $2$ and $K$ -- the exponents are
$mn^2+m(m-1)n^2=(K-1)^4/4$ and $m(K-2)+m(m-1)(K+1)=(K-1)(K^2-7)/4$ -- gives
\begin{equation}\label{eq:global-fourth}
|\Disc(S_K)|
 =2^{(K-1)^4/4}K^{(K-1)(K^2-7)/4}J_K^4
\end{equation}
for the integer
$J_K=\prod_{[r]}I_r\prod_{[r]<[s]}|M_{r,s}|>0$.  Exact resultant
computations through $K=47$ (the normalized prime-to-$2K$ mass
$4\log J_K^{(2K)}/K^4$ increases from $0.0837$ at $K=23$ to $0.1182$ at
$K=47$) show that the residual logarithmic mass genuinely remains
$K^4$-scale on the computed range, so \eqref{eq:global-fourth} is
structural, not a hidden proof of \eqref{eq:sharp-target}.

\section{Fixed-product rows and exact split support}

For one valid subgroup pair put $u=xy$, $s=x+y$.  Then
\begin{equation}\label{eq:rowparam}
c(x)+c(y)=\frac{2(u-1)}{u-s+1}.
\end{equation}
At fixed $u\ne1$, regard the parameter as the finite value
$\sigma=c(x)+c(y)$ in \eqref{eq:rowparam}; the case $u-s+1=0$ is the
excluded infinite parameter.  For finite $\sigma$ -- necessarily
$\sigma\ne0$, since $\sigma=0$ would force $u=1$ -- \eqref{eq:rowparam}
gives $s=u+1-2(u-1)/\sigma$, and $\{x,y\}$ is then the root pair of
$T^2-sT+u$; hence there is at most one unordered pair.
Let $I_K(p)$ be the maximum number of common parameters of two distinct product
rows.  The exact cross-row decomposition gives
\begin{equation}\label{eq:rowenergy}
\Delta_K(p)\le4I_K(p)(K-1)(K-2).
\end{equation}

Normalize a row pair to $(1,r)$ and set
\[
\tau=\zeta+\zeta^{-1},\qquad
\alpha_r=U_{r-1}(\tau/2),\qquad
\beta_r=-2U_{r-2}(\tau/2).
\]
The row equation is the affine trace relation
\begin{equation}\label{eq:affinetrace}
Y=\alpha_rX+\beta_r.
\end{equation}
Let $C_K(T)$ be the minimal polynomial of $\tau$ and set
\[
H_K(T)=(T-2)C_K(T).
\]
The two full row polynomials have one unique characteristic-zero common trace,
the invalid row trace.  Removing it produces fixed-degree source/target
cofactors over $\Z[\tau]$.  For every split embedding prime $\mathfrak q$,
\begin{equation}\label{eq:gcdcount}
\deg\gcd(P_{K,r},Q_{K,r}\bmod\mathfrak q)
 =I_{1,r}(\mathfrak q).
\end{equation}

Let the coefficient ideal of the first three fixed-degree subresultants
contract to $\Lambda_{K,r}^{(3)}\Z$.

\begin{theorem}[Exact real-trace order-three support]\label{thm:lambda}
For every split rational prime $p\equiv1\pmod K$,
\[
p\mid\Lambda_{K,r}^{(3)}
\]
if and only if some row pair in ratio class $[r]$ has at least three common
Cayley parameters.
\end{theorem}

\begin{proof}
Let $S_0,S_1,S_2$ be the first three fixed-degree subresultants of the two
cofactors, and let $J^{\mathrm{tr}}_{K,r}\subset\Z[\tau]$ be the ideal
generated by all their coefficients, so
$\Lambda_{K,r}^{(3)}\Z=J^{\mathrm{tr}}_{K,r}\cap\Z$.  This ideal is nonzero:
the cofactors have no common characteristic-zero trace (the unique common
trace was divided out), so $S_0\ne0$.

We use the contraction criterion: for a nonzero ideal $J$ of $\Z[\tau]$
with $J\cap\Z=\Lambda\Z$, a rational prime $p$ divides $\Lambda$ if and
only if $J\subseteq\mathfrak q$ for some prime ideal $\mathfrak q$ of
$\Z[\tau]$ above $p$.  Indeed, $J\subseteq\mathfrak q$ gives
$\Lambda\in\mathfrak q\cap\Z=p\Z$.  Conversely, if $p\mid\Lambda$ and
$J+p\Z[\tau]$ were the whole ring, localizing at $p$ would put
$1\in J_{(p)}$; clearing denominators in the domain $\Z[\tau]$ then gives
an element of $J\cap\Z$ prime to $p$, contradicting $p\mid\Lambda$.
Hence $J+p\Z[\tau]$ is proper, lies in a maximal ideal $\mathfrak q$, and
$\mathfrak q\cap\Z=p\Z$.

The valid cofactors remain fixed-degree and squarefree after split
reduction -- their leading coefficients are $1$ and the unit
$\alpha_r^{m+1}$, and $p\ne K$ -- so \eqref{eq:gcdcount} is exact and
there is no specialization degree drop.  By the criterion,
$p\mid\Lambda_{K,r}^{(3)}$ if and only if some $\mathfrak q\mid p$
contains every coefficient of $S_0,S_1,S_2$, i.e. if and only if
$S_0\equiv S_1\equiv S_2\equiv0\pmod{\mathfrak q}$.  Brown--Traub's
fixed-degree subresultant criterion \cite{BT1971} identifies this with
$\deg\gcd(P_{K,r},Q_{K,r}\bmod\mathfrak q)\ge3$, and \eqref{eq:gcdcount}
identifies that degree with $I_{1,r}(\mathfrak q)$.  As $\mathfrak q$
ranges over the primes above the split prime $p$, the associated real
embeddings cover all row pairs in ratio class $[r]$, which gives the
rational-prime statement.
\end{proof}

After removing factors $2,K$ and multiplying over ratio classes, write the
result as $\Lambda'_K{}^{(3)}$.  Then
\begin{equation}\label{eq:outside-support}
p\nmid\Lambda'_K{}^{(3)}
\Longrightarrow
I_K(p)\le2,
\quad
\Delta_K(p)\le8(K-1)(K-2),
\quad
E_+(A)\le11(K-1)(K-2).
\end{equation}
The older dehomogenized $D'/G'$ invariants remain valid only as one-way
obstructions because specialization can reduce leading degree at projective
infinity.

\section{Characteristic-zero affine-intercept rigidity}

Write
\[
t_a=\zeta^a+\zeta^{-a},\qquad
s_a=\zeta^a-\zeta^{-a},\qquad
\gamma_r(a,b)=t_b-\alpha_rt_a.
\]
Let $[a]$ denote the trace class $a\sim-a$, and put
$E_j=|1-\zeta^j|$, $1\le j\le m=(K-1)/2$.

\begin{lemma}[Prime-conductor multiplicative independence]\label{lem:bass}
The positive numbers $E_1,\ldots,E_m$ are multiplicatively independent.
\end{lemma}

\begin{proof}
Suppose $\prod_{j=1}^mE_j^{n_j}=1$ with $n_j\in\Z$.  Since
$E_j^2=(1-\zeta^j)(1-\zeta^{-j})$ and
$(1-\zeta^{-j})=-\zeta^{-j}(1-\zeta^j)$, squaring the relation gives
$\prod_j(1-\zeta^j)^{2n_j}=\omega$ for a root of unity $\omega$ (a sign
times a power of $\zeta$).  Taking norms from $\mathbf Q(\zeta)$: every
$1-\zeta^j$ is a Galois conjugate of $1-\zeta$, so
$N(1-\zeta^j)=N(1-\zeta)=\Phi_K(1)=K$, while $N(\omega)=1$.  Hence
$K^{2\sum_jn_j}=1$, forcing
\[
\sum_{j=1}^mn_j=0.
\]
The relation therefore rewrites as one among the ratios
$(1-\zeta^j)/(1-\zeta)$, $2\le j\le m$ (the $j=1$ ratio is $1$, and the
vanishing exponent sum absorbs the powers of $1-\zeta$).  By Bass's
theorem for prime conductor \cite{Bass1966} (Theorem~3), these ratios are
multiplicatively independent modulo torsion: their classes form a free
basis of the cyclotomic units modulo torsion.  Hence $2n_j=0$ for
$j\ge2$, and the vanishing sum gives $n_1=0$.
\end{proof}

Use the elementary identities
\begin{equation}\label{eq:absfacts}
|s_n|=E_{[2n]},\qquad
|t_x-t_y|=E_{[x-y]}E_{[x+y]}.
\end{equation}

\begin{theorem}[Trace-intercept fibre rigidity]\label{thm:intercept}
For every odd prime $K\ge7$ and every normalized ratio $r\ne1$, every fibre of
$([a],[b])\mapsto\gamma_r(a,b)$ has cardinality at most two.

If $r=-1$, a double fibre is only a swap.  If $r\ne\pm1$, every double fibre is,
up to signs and exchange,
\begin{equation}\label{eq:doublefibre}
([a],[a+1-r]),\qquad([a+2],[a+1+r]).
\end{equation}
\end{theorem}

\begin{proof}
Suppose two distinct pairs have the same intercept.  For $r\ne-1$,
\[
t_b-t_d=\alpha_r(t_a-t_c).
\]
Taking absolute values and applying \eqref{eq:absfacts} yields
\[
E_{[b-d]}E_{[b+d]}E_{[2]}
 =E_{[a-c]}E_{[a+c]}E_{[2r]}.
\]
Lemma~\ref{lem:bass} makes the multisets of folded indices equal.  Since
$[2r]\ne[2]$, one source difference must be $[2]$; after changing signs and
ordering take $c=a+2$.  The remaining two indices force the target classes in
\eqref{eq:doublefibre}.  The equality itself follows from
\[
t_{x+r}-t_{x-r}=s_rs_x,
\qquad
t_{a+2}-t_a=s_1s_{a+1}.
\]
A third point would form a triangle in the step-two quotient of the odd
$K$-cycle by $a\sim-a$; that quotient is a path.  For $r=-1$ the intercept is
$t_a+t_b$.  If $t_a+t_b=t_c+t_d$, subtracting gives a rational relation
among the powers $1,\zeta,\ldots,\zeta^{K-1}$ whose coefficient sum is
$4-4=0$.  The space of such relations is spanned by $(1,\ldots,1)$, the
coefficient vector of $\Phi_K$; the zero coefficient sum forces the scalar
to vanish, so every coefficient vanishes separately, the multisets
$\{\pm a,\pm b\}$ and $\{\pm c,\pm d\}$ in $\mathbf F_K$ agree, and
$\{[a],[b]\}=\{[c],[d]\}$: only identity or swap.  This covers in
particular $K=7$, where the relation has length $8>\deg\Phi_K=6$.
\end{proof}

The distinguished characteristic-zero row solution $([1],[r])$ is therefore
isolated.  This mechanism is related to cyclotomic-unit cross-ratio methods in
projective equivalence of roots of unity \cite{Fu2022}, but the affine theorem
above is proved directly in the prime-conductor setting.
